% !TEX root = ../Main.tex
\chapter{力的合成：三种特殊三角形}

两个力 $\vec{F_1}, \vec{F_2}$ 的合力由\textbf{平行四边形法则}给出，大小由余弦定理：
\[
|\vec F_1 + \vec F_2|^2 = F_1^2 + F_2^2 + 2 F_1 F_2 \cos \alpha,
\]
其中 $\alpha$ 为两力的\textbf{夹角}。但\textbf{三种特殊夹角}（$60^\circ, 90^\circ, 120^\circ$）下，等大情况 $F_1 = F_2 = F$ 能得到\textbf{记忆级}的合力：

\section{三种特殊合力}

\state{等大两力的三种特殊夹角}{
    设 $F_1 = F_2 = F$，两力夹角为 $\alpha$，则合力大小
    \[
    F_{\text{合}} = 2 F \cos \tfrac{\alpha}{2}.
    \]
    代入三种特殊值：
    \begin{center}
    \begin{tabular}{c|c|c}
        \hline
        夹角 $\alpha$ & 合力 & 合成三角形 \\
        \hline
        $60^\circ$ & $\sqrt 3 F$ & 等腰三角形，顶角 $120^\circ$ \\
        $90^\circ$ & $\sqrt 2 F$ & 等腰直角三角形 \\
        $120^\circ$ & $F$ & 等边三角形 \\
        \hline
    \end{tabular}
    \end{center}
}

\begin{center}
\begin{tikzpicture}[>=Stealth, scale=0.9]
    % 60° case
    \draw[->, very thick, blue] (0, 0) -- (1.5, 0) node[right] {\small $F_1$};
    \draw[->, very thick, blue] (0, 0) -- ({1.5*cos(60)}, {1.5*sin(60)}) node[above left] {\small $F_2$};
    \draw[->, very thick, red] (0, 0) -- ({1.5+1.5*cos(60)}, {1.5*sin(60)}) node[right] {\small $\sqrt 3 F$};
    \draw[dashed] (1.5, 0) -- ({1.5+1.5*cos(60)}, {1.5*sin(60)});
    \draw[dashed] ({1.5*cos(60)}, {1.5*sin(60)}) -- ({1.5+1.5*cos(60)}, {1.5*sin(60)});
    \node at (1.2, -0.6) {\small $\alpha = 60^\circ$};
    % 90° case
    \begin{scope}[xshift=5cm]
        \draw[->, very thick, blue] (0, 0) -- (1.5, 0) node[right] {\small $F_1$};
        \draw[->, very thick, blue] (0, 0) -- (0, 1.5) node[above] {\small $F_2$};
        \draw[->, very thick, red] (0, 0) -- (1.5, 1.5) node[above right] {\small $\sqrt 2 F$};
        \draw[dashed] (1.5, 0) -- (1.5, 1.5);
        \draw[dashed] (0, 1.5) -- (1.5, 1.5);
        \node at (0.8, -0.6) {\small $\alpha = 90^\circ$};
    \end{scope}
    % 120° case
    \begin{scope}[xshift=10cm]
        \draw[->, very thick, blue] (0, 0) -- (1.5, 0) node[right] {\small $F_1$};
        \draw[->, very thick, blue] (0, 0) -- ({1.5*cos(120)}, {1.5*sin(120)}) node[above left] {\small $F_2$};
        \draw[->, very thick, red] (0, 0) -- ({1.5+1.5*cos(120)}, {1.5*sin(120)}) node[above right] {\small $F$};
        \draw[dashed] (1.5, 0) -- ({1.5+1.5*cos(120)}, {1.5*sin(120)});
        \draw[dashed] ({1.5*cos(120)}, {1.5*sin(120)}) -- ({1.5+1.5*cos(120)}, {1.5*sin(120)});
        \node at (0.5, -0.6) {\small $\alpha = 120^\circ$};
    \end{scope}
\end{tikzpicture}
\end{center}

\rmkb{
    \textbf{记忆要点}：
    \begin{itemize}
        \item \textbf{$60^\circ + 60^\circ = 120^\circ$}（两力夹角 $60^\circ$）——合力 $\sqrt 3 F$，合成三角形中顶角 $120^\circ$；
        \item \textbf{$120^\circ + 120^\circ + 120^\circ$}（两力夹角 $120^\circ$）——三角形为等边，\textbf{合力 $= F$}；
        \item \textbf{$45^\circ + 45^\circ = 90^\circ$}（两力夹角 $90^\circ$）——合力 $\sqrt 2 F$。
    \end{itemize}
    \textbf{考试经验}：看到 $60^\circ, 90^\circ, 120^\circ$ 等\textbf{明显的特殊角}时，别套余弦定理浪费时间——直接写合力值。
}

\section{一般的合成公式}

\state{任意两力合成}{
    对任意 $F_1, F_2$ 及夹角 $\alpha$：
    \[
    F_{\text{合}} = \sqrt{F_1^2 + F_2^2 + 2 F_1 F_2 \cos \alpha}.
    \]
    \textbf{特殊情形}：
    \begin{itemize}
        \item $\alpha = 0^\circ$：$F_{\text{合}} = F_1 + F_2$（同方向相加）；
        \item $\alpha = 180^\circ$：$F_{\text{合}} = |F_1 - F_2|$（反方向相减）；
        \item $\alpha = 90^\circ$：$F_{\text{合}} = \sqrt{F_1^2 + F_2^2}$（直角和）。
    \end{itemize}

    \textbf{范围}：任意 $\alpha$ 下，
    \[
    |F_1 - F_2| \le F_{\text{合}} \le F_1 + F_2.
    \]
}

\ex[三力平衡的特殊角]{
    三个共点力 $F_1 = F_2 = F_3 = F$，相互夹角均为 $120^\circ$。证明合力为零。
}

\sol{
    取两力合成：$F_1 + F_2$（夹角 $120^\circ$，等大）的合力为 $F$，方向与 $F_3$ \textbf{反向}。故三力合力 $= F - F = 0$。
    \textbf{这是"三力共点成 $120^\circ$ 等大 $\Rightarrow$ 平衡"的速记：等边三角形闭合 $\Rightarrow$ 零合力。}
}
